% Source-aligned English for printed p.173.
\SourcePage{173}

\LayerTarget{tohoku:proof:3.7.1:2:p2}
\noindent inducing $g_1$. The support of $h$ is contained in
$A\cap f^{-1}(B_1)$, and hence belongs to $\Psi'$. Thus $h$ may be regarded
as an element of $\Gamma_\Phi(f_\Psi(F))$, evidently inducing the same section
as $g'$ on $U_1$, and consequently inducing $g$ on $B$.

\LayerTarget{tohoku:construction:relative-leray-spectral-functor}
It follows from the corollary to Proposition 3.3.1 that, under either set of
conditions in Lemma 3.7.1, Theorem 2.4.1 may be applied to the composite
functor $\Gamma_{\Psi'}=\Gamma_\Phi f_\Psi$. There is a cohomological spectral
functor on $\mathbf C_Y$ abutting to the graded functor $(H^n_{\Psi'}(Y,F))$, whose
initial term is
\LayerTarget{tohoku:eq:relative-leray-preliminary-E2}
\[
E_2^{p,q}=H^p_\Phi\bigl(X,(R^qf_\Psi)(F)\bigr).
\]
It remains to describe the sheaves $(R^qf_\Psi)(F)$. More generally:

\LayerTarget{tohoku:lemma:3.7.2}
\textsc{Lemma 3.7.2.} \textit{Let $T$ be a covariant functor
from an abelian category $\mathbf C$ to the category $\mathbf C^X$ of sheaves of abelian
groups on $X$. Suppose that every object of $\mathbf C$ is isomorphic to a subobject
of an injective object, so that the right derived functors $R^qT$ exist. Then,
for every $A\in\mathbf C$, the sheaf $R^qT(A)$ is identified with the sheaf associated
\textup{(}cf. \textup{3.1)} to the presheaf that assigns $R^q(\Gamma^U T)(A)$ to every open
subset $U\subset X$, where $\Gamma^U$ denotes the functor
$F\to\Gamma(U,F)$ on $\mathbf C^X$.}

\LayerTarget{tohoku:proof:3.7.2}
Indeed, let $C(A)$ be the complex associated with a right resolution of $A$ by
injective objects. Then $R^qT(A)=H^q(T(C(A)))$. The $q$-th cohomology sheaf of
a complex of sheaves $K=(K^i)$ is precisely the sheaf associated to the
presheaf that assigns the group $H^q(\Gamma(U,K))$ to an open subset $U$.
Thus $R^qT(A)$ is the sheaf associated to the presheaf
$U\to H^q\bigl(\Gamma(U,T(C(A)))\bigr)
=R^q(\Gamma^U T)(A)$,
which proves the lemma.

\LayerTarget{tohoku:eq:derived-direct-image-presheaf}
In the present case $T=f_\Psi$, it follows that $R^qf_\Psi(F)$ is the sheaf
associated to the presheaf
$U\to R^q(\Gamma^U f_\Psi)(F)=R^q\Gamma_{\Psi(U)}(F)$.
We have already observed, as an immediate consequence of Proposition 3.1.3,
that the derived functors of
$\Gamma_{\Psi(U)}(f^{-1}(U),F)$ are precisely the
$H^q_{\Psi(U)}(f^{-1}(U),F)$. We therefore obtain:

\LayerTarget{tohoku:theorem:3.7.3}
\textsc{Theorem 3.7.3.} \textit{Let $f$ be a continuous map from
a space $Y$ to a space $X$, and suppose that $X$ and $Y$ are endowed with
antifilters of closed subsets $\Phi$ and $\Psi$, respectively. The notation
$\Psi(U)$, $f_\Psi$, and $\Psi'$ is as at the beginning of this section.
Suppose that $\Phi$ is paracompactifying or that $\Psi$ is the set of all
closed subsets of $Y$. Then there exists a cohomological spectral functor on
the category $\mathbf C^Y$ of sheaves of abelian groups $F$ on $Y$, abutting to the
graded functor $(H^n_{\Psi'}(Y,F))$, whose initial term is}
\LayerTarget{tohoku:eq:relative-leray-E2}
\[
E_2^{p,q}(F)=H^p_\Phi(X,R^qf_\Psi(F)).
\]
\LayerTarget{tohoku:def:relative-derived-direct-image}
\textit{Here $R^qf_\Psi(F)$ is the sheaf associated to the presheaf on $X$
that sends an open subset $U$ to $H^q_{\Psi(U)}(f^{-1}(U),F)$.}

\LayerTarget{tohoku:specialization:ordinary-leray}
The simplest case is that in which $\Phi$ and $\Psi$ are the sets of all
closed subsets of $X$ and $Y$, respectively. Then, \textit{without any hypothesis} on $X,Y,f$, \textit{one obtains a spectral functor abutting to
$(H^n(Y,F))$, whose initial term is
\LayerTarget{tohoku:eq:ordinary-leray-E2}
$H^p(X,R^qf(F))$, where $R^qf(F)$ is the sheaf associated to the presheaf
$U\to H^q(f^{-1}(U),F)$.}
\LayerTarget{tohoku:application:leray-algebraic-varieties:p1}
This statement may, for example, be applied usefully in the cohomology theory

\TohokuContinuation{174}
